\documentclass{article}
\usepackage{amsmath,amssymb}
\usepackage[margin=1in]{geometry}
\allowdisplaybreaks
\usepackage[UTF8]{ctex}
\begin{document}

\section*{单车辆类型的ADP车队管理政策}

对于每一个时间段$t\in\mathcal T$，在当前的状态$r_t$、观测到的负载情况$d_t$，以及提供的策略$\pi$下，求解相应的问题。
\begin{equation}
\min_{(x_t,r_{t+1},z_{t+1})}\;
\sum_{i,j\in\mathcal I}\sum_{l\in\mathcal L}c_{ijlt}x_{ijlt}
+\sum_{j\in\mathcal I}\sum_{q=1}^{R}\tilde v^\pi_{j,t+1}(q)z_{j,t+1}(q).
\tag{11}
\end{equation}
在以下条件下：
\begin{align}
\sum_{j\in\mathcal I}\sum_{l\in\mathcal L}x_{ijlt}
&=r_{it} &&\text{for all }i\in\mathcal I,\tag{12}\\
\sum_{i\in\mathcal I}\sum_{l\in\mathcal L}x_{ijlt}-r_{j,t+1}
&=0 &&\text{for all }j\in\mathcal I,\tag{13}\\
r_{j,t+1}-\sum_{q=1}^{R}z_{j,t+1}(q)
&=0 &&\text{for all }j\in\mathcal I,\tag{14}\\
x_{ijlt}&\leq d_{ijlt} &&\text{for all }i,j\in\mathcal I,\ l\in\mathcal L,\tag{15}\\
z_{j,t+1}(q)&\leq 1 &&\text{for all }j\in\mathcal I,\ q=1,\ldots,R,\tag{16}\\
x_{ijlt},\ r_{j,t+1},\ z_{j,t+1}(q)&\in\mathbb Z_+
&&\text{for all }i,j\in\mathcal I,\ l\in\mathcal L,\ q=1,\ldots,R.
\tag{17}
\end{align}

\section*{范围、数据以及变量角色}

规划范围、位置、车辆类型以及移动模式分别如下：
\[
\mathcal T=\{1,\ldots,T\},\qquad
\mathcal I,\qquad
\mathcal K\ \text{with }|\mathcal K|=1,\qquad
\mathcal L=\{0,\ldots,L\}.
\]
如论文第2节所述，此公式中的车辆类型上标部分被省略了。$R$表示可用的车辆总数；它与下面的状态转移函数$R^\pi_{t+1}$是不同的。移动模式$0$表示空置状态的重新定位；其余的模式则表示承载不同类型的货物。

决策$x_{ijlt}$指的是在时间$t$，使用移动模式$l$，从位置$i$派遣到位置$j$的车辆数量。固定系数$c_{ijlt}$表示派遣一辆这样的车辆所需的成本。状态组件$r_{it}$则表示在时间$t$，位置$i$处可用的车辆数量。因此，在周期$t$问题中，$r_t$是已知的，而$r_{t+1}$则由该问题决定。变量$z_{j,t+1}(q)$包含了所提供的分段线性连续函数；它们并不是额外的车辆类型或负载类型。

$D_{ijlt}$表示在时间$t$，以移动模式$l$下，需要从$i$运输到$j$的随机数量的负载。而$d_{ijlt}$则代表一个特定的实现情况。省略索引表示该向量包含的是这些索引所代表的各个组成部分；具体来说，
\[
\begin{aligned}
r_t&=\{r_{it}:i\in\mathcal I\},\\
d_t&=\{d_{ijlt}:i,j\in\mathcal I,\ l\in\mathcal L\},\\
d&=\{d_{ijlt}:i,j\in\mathcal I,\ l\in\mathcal L,\ t\in\mathcal T\}
  =\{d_t:t\in\mathcal T\}.
\end{aligned}
\]
空置的重新定位动作在负载可用性的方面是没有限制的：
\[
D_{ij0t}=\infty
\qquad\text{for all }i,j\in\mathcal I,\ t\in\mathcal T.
\]
如果仅使用该车辆类型来运输负载类型$l$，则无法实现该运输任务。因此，本文通过以下方式对此进行了处理：
\[
c_{ijlt}=\infty
\qquad\text{for all }i,j\in\mathcal I,\ t\in\mathcal T
\quad\text{for that load type }l.
\]
论文中提到，通常情况下，$c_{ijlt}<0$代表$l\in\mathcal L\setminus\{0\}$；但这并不是一种普遍适用的输入限制。

在任何两个地点之间旅行都需要一个时间周期。关于未来负载的信息是未知的。在当前周期中未涵盖的负载将由紧急分包商来处理，然后立即离开系统。该模型采用成本最小化方法来实现预期利润的最大化，这一方法适用于有限的时间范围。第1节中提出的随机假设是：给定$r_t$，$D_t$与$\{D_{t'}:t'=1,\ldots,t-1\}$是独立的。论文还假设$c_t$会受到微小的随机扰动，从而得到一个最优解，并且决策和转换函数被正确定义；不过，论文并未具体说明扰动的大小或分布情况。

\section*{请提供后续的函数及其嵌入方式。}

政策$\pi$是固定不变的，并未在(11)中经过优化处理。其后续近似解是可以分离的。对于单一车辆类型而言，第7页中所述的(9)的特化情况，是可以实现的。
\begin{equation}
\tilde V^\pi_t(r_t)=\sum_{i\in\mathcal I}\tilde V^\pi_{it}(r_{it}).
\tag{9, single type}
\end{equation}
每个组成部分都是一个一维的、分段线性凸函数，其不可微点集是正整数的一个子集。$\tilde V^\pi_{it}(\cdot)$的相关定义域为
\[
\{0,1,\ldots,R\}.
\]
所选的表示方式采用了纸张的标准化处理和坡度差异处理：
\[
\tilde V^\pi_{it}(0)=0,\qquad
\tilde v^\pi_{it}(q)=\tilde V^\pi_{it}(q)-\tilde V^\pi_{it}(q-1),
\qquad q=1,\ldots,R.
\]
这里的$\tilde v^\pi_{it}(q)$表示的是开区间$(q-1,q)$上的斜率。这些数字都是给定的系数，而不是阶段决策的结果。对于(11)中出现的连续分量来说，凸性起到了重要作用。
\[
\tilde v^\pi_{j,t+1}(1)\leq\tilde v^\pi_{j,t+1}(2)
\leq\cdots\leq\tilde v^\pi_{j,t+1}(R),
\qquad j\in\mathcal I.
\]
如第8页所述，在(11)--(17)的最优解中，
\[
\sum_{q=1}^{R}\tilde v^\pi_{j,t+1}(q)z_{j,t+1}(q)
=\sum_{q=1}^{r_{j,t+1}}\tilde v^\pi_{j,t+1}(q)
=\tilde V^\pi_{j,t+1}(r_{j,t+1}),
\qquad j\in\mathcal I.
\]
因此，(11)中的第二项的计算结果为：
\[
\sum_{j\in\mathcal I}\tilde V^\pi_{j,t+1}(r_{j,t+1}).
\]
这些就是来源中的最优解恒等式，它们解释了嵌入过程。对于$z$来说，并没有额外的排序约束。约束条件（13）和（14）则分别保留下来，就像在选定的打印版本中那样。

\section*{继承的可行集合与策略定义}

用于定义策略的可行集合，包括调度决策以及它们所产生的下一期状态。将定义行（1）至（4）替换为$|\mathcal K|=1$，并去掉类型上标后，即可得到完整的结果。
\begin{equation}
\mathcal Y(r_t,d_t)=
\left\{(x_t,r_{t+1}):
\begin{aligned}
&\sum_{j\in\mathcal I}\sum_{l\in\mathcal L}x_{ijlt}=r_{it}
&&\text{for all }i\in\mathcal I,\\
&\sum_{i\in\mathcal I}\sum_{l\in\mathcal L}x_{ijlt}-r_{j,t+1}=0
&&\text{for all }j\in\mathcal I,\\
&x_{ijlt}\leq d_{ijlt}
&&\text{for all }i,j\in\mathcal I,\ l\in\mathcal L,\\
&x_{ijlt}\in\mathbb Z_+
&&\text{for all }i,j\in\mathcal I,\ l\in\mathcal L
\end{aligned}
\right\}.
\tag{1--4, single type}
\end{equation}
这些就是那些已经体现在(12)、(13)、(15)和(17)中的物理约束条件，而不是需要额外解决的新的公式。

马尔可夫确定性策略的决策函数为$X_t^\pi(\cdot,\cdot)$，它将$(r_t,d_t)$映射到$x_t$；而转移函数则为$R_{t+1}^\pi(\cdot,\cdot)$，它将$(r_t,d_t)$映射到$r_{t+1}$，对于$t\in\mathcal T$则如此。对于所提供的近似值，其定义方程为：
\begin{equation}
\bigl(X_t^\pi(r_t,d_t),R_{t+1}^\pi(r_t,d_t)\bigr)
=\operatorname*{argmin}_{(x_t,r_{t+1})\in\mathcal Y(r_t,d_t)}
\left\{c_t\cdot x_t+\tilde V^\pi_{t+1}(r_{t+1})\right\}.
\tag{10}
\end{equation}
方程(11)以及(12)--(17)都是该论文中明确表述的相同政策定义。精确的最优值函数方程(7)--(8)被(10)中的近似值所替代；这些近似值并不适用于$\pi$。

\section*{在视野范围内进行的操作，以及累积成本}

给定初始状态$r_1$和加载轨迹$d$，该策略生成的状态和决策轨迹满足以下条件：
\begin{equation}
x_t^{\pi d}=X_t^\pi(r_t^{\pi d},d_t),\qquad
r_{t+1}^{\pi d}=R_{t+1}^\pi(r_t^{\pi d},d_t),\qquad
r_1^{\pi d}=r_1,
\quad t\in\mathcal T.
\tag{18}
\end{equation}
该政策下的实际累计成本是通过递归方式定义的。
\begin{equation}
\begin{aligned}
F_t^\pi(r_t,d_t,d_{t+1},\ldots,d_T)
={}&c_t\cdot X_t^\pi(r_t,d_t)\\
&+F_{t+1}^\pi\bigl(R_{t+1}^\pi(r_t,d_t),d_{t+1},d_{t+2},\ldots,d_T\bigr),
\end{aligned}
\qquad t\in\mathcal T,
\tag{5}
\end{equation}
带有打印出的边界条件的
\[
F_{T+1}^\pi(\cdot)=0.
\]
重复应用后，会得到打印出的地平线扩展效果。
\begin{equation}
\begin{aligned}
F_1^\pi(r_1,d_1,\ldots,d_T)
={}&c_1\cdot X_1^\pi(r_1,d_1)
+c_2\cdot X_2^\pi\bigl(R_2^\pi(r_1,d_1),d_2\bigr)+\cdots\\
&+c_T\cdot X_T^\pi\bigl(R_T^\pi(R_{T-1}^\pi(\ldots,d_{T-2}),d_{T-1}),d_T\bigr).
\end{aligned}
\tag{6}
\end{equation}
相应的预期总成本为
\[
\mathbb E\bigl\{F_1^\pi(r_1,D_1,\ldots,D_T)\mid r_1\bigr\}.
\]
累积成本$F_t^\pi(\cdot,d_t,\ldots,d_T)$并不具有可分离性。特别是，该论文并未提出这种身份关系。
\[
\tilde V_t^\pi(r_t)
=\mathbb E\bigl\{F_t^\pi(r_t,D_t,\ldots,D_T)\mid r_t\bigr\}
\]
对于这种近似策略而言，所显示的等式并不构成操作上的约束条件。

\paragraph{源范围与终端遗漏。}
上述表述采用了所提供的原始论文，即第4页至第8页的内容，其中包括第6页的文字以及第5页、第7页和第8页的图片。该论文将近似函数描述为$\{\tilde V_t^\pi(\cdot):t\in\mathcal T\}$，但(10)--(11)部分则使用了$t+1$的后续内容，而$T+1$则包含在$t=T$中。所提供的论文并未给出$\tilde V_{T+1}^\pi$或其斜率的终止值。对于$F_{T+1}^\pi$的打印边界条件，仅保留了$F^\pi$的情况；没有添加任何零终止近似或其他修正措施。所引用的基于抽样的构建方法来自Godfrey和Powell (2002a)，但并未提供相关文献，因此此处不予验证；该方法并非所选固定策略表述的一部分。

\end{document}
